\chapter{Usher Syndrome}
\index{Usher Syndrome}
\label{sec:Usher Syndrome}

\section{Overview}
\begin{itemize}[noitemsep]
  \item Usher syndrome is a rare inherited condition causing both hearing loss and vision loss from retinitis pigmentosa
  \item It is the most common cause of combined deafness and blindness (deaf-blindness) worldwide
  \item It affects about 1 in 25,000 people and is inherited in an autosomal recessive pattern
\end{itemize}
\section{Causes}
This condition happens because of mutations in multiple genes (USH1 and USH2 families, including MYO7A, USH2A, CDH23) that are inherited in an autosomal recessive pattern. The affected genes encode proteins important for the function of sensory cells in both the inner ear and the retina, so defects cause progressive retinal degeneration and inner-ear dysfunction. Balance problems can also occur.
\section{Risk Factors}
\begin{itemize}[noitemsep]
  \item Family history of Usher syndrome
  \item Having two carrier parents (autosomal recessive inheritance)
  \item Consanguinity increases the risk
  \item Certain populations have higher frequencies of specific mutations
\end{itemize}
\section{Signs and Symptoms}
\subsection{Common symptoms}
\begin{itemize}[noitemsep]
  \item Hearing loss ranging from profound at birth (type 1) to moderate progressive loss in adolescence (type 2)
  \item Night blindness beginning in childhood or adolescence
  \item Progressive loss of peripheral vision (tunnel vision)
  \item Central vision loss later in life
  \item Balance problems (vestibular dysfunction), especially in type 1
\end{itemize}
\subsection{Emergency warning signs}
\begin{itemize}[noitemsep]
  \item Rapidly worsening vision or sudden hearing loss requires prompt evaluation
\end{itemize}
\section{Progression and Stages}
Usher syndrome is classified into three clinical types based on severity and onset of hearing loss. Type 1: profound congenital deafness, balance problems, and night blindness appearing in early childhood. Type 2: moderate hearing loss present at birth, normal balance, and retinitis pigmentosa developing in adolescence. Type 3: progressive hearing and vision loss beginning in adulthood. Vision loss is progressive in all types.
\section{Complications}
\begin{itemize}[noitemsep]
  \item Combined deafness and blindness (deaf-blindness)
  \item Severe visual field constriction and eventual blindness
  \item Difficulty with communication, mobility, and daily activities
  \item Social isolation and psychological distress
  \item Increased risk of accidents and falls
  \item Barriers to education and employment
\end{itemize}
\section{Diagnosis}
\begin{itemize}[noitemsep]
  \item Audiological testing to document hearing loss
  \item Dilated fundus examination showing retinitis pigmentosa changes
  \item Electroretinography (ERG) to assess retinal function
  \item Visual field testing
  \item Vestibular testing when balance problems are present
  \item Genetic testing to confirm the diagnosis and type
\end{itemize}
\section{Prevention}
\begin{itemize}[noitemsep]
  \item Genetic counseling for affected families
  \item Carrier screening and prenatal diagnosis where family mutations are known
  \item Newborn hearing screening allows early detection and intervention
  \item Avoidance of smoking, which may accelerate retinal degeneration
\end{itemize}
\section{Treatment Options}
\begin{itemize}[noitemsep]
  \item There is no cure, and treatment focuses on supporting both senses
  \item Hearing aids and cochlear implants for hearing loss
  \item Speech and language therapy and sign language training
  \item Low-vision rehabilitation and mobility (orientation) training
  \item Gene therapy and clinical trials are under investigation for specific Usher genes
  \item Vitamin A palmitate may slow retinal degeneration in some patients (under specialist guidance)
  \item Assistive technology for communication and daily living
\end{itemize}
\section{Living with It}
\begin{itemize}[noitemsep]
  \item Begin hearing rehabilitation and visual rehabilitation early
  \item Use a coordinated deaf-blind services team
  \item Use assistive and accessible technology (captioning, screen readers, tactile aids)
  \item Join deaf-blind support communities and advocacy organizations
  \item Stay informed about research and gene-therapy advances
\end{itemize}
\section{Outlook}
Usher syndrome is lifelong and progressive, with vision and hearing declining over decades. With early intervention, hearing aids or implants, and comprehensive rehabilitation, most people maintain communication, mobility, and independence. Research into gene therapy offers future hope for slowing the condition.
